\chapter{Mesure, évolution unitaire et dynamique quantique}

\lectureoverview{Nous rassemblons d’abord les principes de mesure déjà établis : les observables sont des opérateurs hermitiens, leurs valeurs propres sont les résultats possibles, les états distinguables sont orthogonaux et la règle de Born attribue les probabilités. La nouvelle question porte sur la dynamique. La réversibilité exige une évolution linéaire commune qui préserve les produits scalaires ; l’opérateur d’évolution est donc unitaire et son générateur infinitésimal, hermitien. Ce générateur est le Hamiltonien, dont l’action donne l’équation de Schr\"odinger dépendant du temps. Nous déduisons ensuite la formule de la valeur moyenne, l’équation des commutateurs pour les moyennes, son lien avec les crochets de Poisson classiques et la conservation de l’énergie moyenne.}

Pour l’instant, laissons de côté le système de spin. Nous y reviendrons une fois la structure générale établie. Les principes qui suivent ne constituent pas un ensemble minimal et indépendant. Certains se renforcent mutuellement ou impliquent en partie les autres, mais les énoncer séparément permet de mieux faire ressortir leur contenu physique.

\section{Ce qu’une mesure peut donner}

Une \emph{observable} est une grandeur qui peut être mesurée. Le terme consacré est un peu plus pompeux que nécessaire ; \emph{mesurable} aurait tout aussi bien convenu. À chaque observable, nous associons un opérateur hermitien. Notons provisoirement cet opérateur \(L\) :
\begin{equation}
L^\dagger=L.
\end{equation}
Tout opérateur linéaire n’est pas hermitien. Inversement, savoir si un opérateur hermitien donné peut être mesuré avec le matériel et les appareils disponibles peut être une question expérimentale difficile. La règle de travail abstraite reste néanmoins généreuse : un opérateur hermitien représente une observable possible, même si la mesure pratique n’a pas encore été inventée.

Si
\begin{equation}
L\ket{\lambda}=\lambda\ket{\lambda},
\end{equation}
alors \(\lambda\) est un résultat numérique possible de la mesure de \(L\). Nous utilisons le même symbole pour désigner la valeur propre et son vecteur propre. Les crochets du ket distinguent les deux : \(\lambda\) est un nombre, tandis que \(\ket{\lambda}\) est un état.

Pourquoi les opérateurs hermitiens ? Un appareil rudimentaire donne des nombres réels. Si un instrument semble fournir un nombre complexe, il donne en réalité deux grandeurs réelles indépendantes. Les valeurs propres d’une observable doivent donc être réelles. En outre, les états associés à des résultats définis différents doivent être physiquement distinguables et, par conséquent, orthogonaux. Dans le cadre de dimension finie utilisé ici, un opérateur linéaire possédant une base propre orthonormée complète et des valeurs propres réelles est hermitien. Les énoncés physiques et la structure hermitienne concordent donc.

\subsection{La distinguabilité équivaut à l’orthogonalité}

L’expression \emph{physiquement distinguables} a un sens opérationnel. Supposons qu’un spin soit préparé dans l’état haut ou dans l’état bas. Une mesure suivant l’axe \(z\) distingue sans ambiguïté les deux possibilités : l’état haut donne \(+1\), l’état bas \(-1\). De même, une mesure suivant l’axe \(x\) permet de distinguer la gauche de la droite.

Supposons maintenant que la préparation soit l’état haut ou l’état droit. Une mesure selon \(z\) donne certainement \(+1\) pour l’état haut, mais l’état droit donne aussi \(+1\) avec une probabilité non nulle. Une mesure selon \(x\) ne fait qu’inverser l’asymétrie. Incliner l’appareil de \(45^\circ\) ne supprime pas le recouvrement entre les distributions des résultats. Aucune orientation ne permet d’obtenir un résultat qui distingue toujours les deux préparations.

\begin{definition}
Deux états sont physiquement distinguables lorsqu’une expérience, ou un ensemble d’expériences, permet de déterminer sans ambiguïté laquelle des préparations a été effectuée.
\end{definition}

La mécanique quantique représente géométriquement la distinguabilité parfaite :
\begin{equation}
\text{distinguishable states}
\quad\Longrightarrow\quad
\text{orthogonal vectors}.
\end{equation}

\begin{classroomqa}
\audiencequestion{Un autre axe de mesure permettrait-il de distinguer un état haut d’un état droit ?}
\lecturerresponse{Non. Une mesure selon \(z\) ou \(x\) peut parfois identifier la préparation, mais au moins un résultat reste compatible avec les deux. Un axe intermédiaire laisse lui aussi des probabilités qui se recouvrent. L’obstacle ne tient pas à un mauvais choix d’appareil : les deux vecteurs d’état ne sont pas orthogonaux.}
\end{classroomqa}

\section{La règle de Born}

Les valeurs propres nous indiquent ce qui peut se produire, mais pas à quelle fréquence. La règle de probabilité fournit le contenu prédictif. Si un système normalisé est préparé dans l’état \(\ket{A}\) et si un vecteur propre normalisé \(\ket{\lambda}\) appartient à l’observable mesurée \(L\), alors
\begin{equation}
P(\lambda)=\left|\braket{\lambda}{A}\right|^2.
\label{eq:l4-born-rule}
\end{equation}
Le recouvrement \(\braket{\lambda}{A}\) est une \emph{amplitude de probabilité}. Il est généralement complexe et ne peut donc pas être lui-même une probabilité. Son carré du module est réel et positif ou nul.

\begin{figure}[htbp]
\centering
\includegraphics[width=0.86\textwidth]{lecture_04_figure_01.png}
\caption{La règle de Born, présentée à côté des principes de mesure : la probabilité est le carré du module d’un recouvrement.}
\lectureframe{00:12:40}
\end{figure}

Choisissons une base propre orthonormée de \(L\) et développons
\begin{equation}
\ket{A}=\sum_i \alpha_i\ket{\lambda_i},
\qquad
\alpha_i=\braket{\lambda_i}{A}.
\end{equation}
La règle de Born devient
\begin{equation}
P(\lambda_i)=|\alpha_i|^2.
\end{equation}
La normalisation donne alors
\begin{equation}
\sum_i P(\lambda_i)=\sum_i|\alpha_i|^2=1.
\end{equation}
La positivité et la normalisation ne suffisent pas à elles seules à déduire la règle de Born, mais elles montrent pourquoi le carré du module s’accorde si naturellement avec la théorie des probabilités. En définitive, cette règle est vérifiée par l’expérience.

\begin{classroomqa}
\audiencequestion{Le carré du module est-il également nécessaire pour que la somme de toutes les probabilités soit égale à un ?}
\lecturerresponse{Pour un état normalisé développé dans une base orthonormée, la somme des carrés des amplitudes vaut bien un. La positivité et la normalisation concordent donc parfaitement avec la règle de Born. Il importe moins ici de savoir lequel de ces énoncés doit être considéré comme logiquement premier que de constater leur cohérence mutuelle ; c’est l’expérience qui fixe la règle de probabilité.}
\end{classroomqa}

\begin{classroomqa}
\audiencequestion{Pourquoi une valeur propre, qui est un nombre, est-elle écrite dans un ket ?}
\lecturerresponse{La notation \(\ket{\lambda}\) désigne le vecteur propre associé à la valeur propre numérique \(\lambda\). Une notation plus pointilleuse pourrait introduire un autre indice, mais trop d’indices finissent vite par obscurcir le calcul. Le symbole ket permet de distinguer le nombre du vecteur.}
\end{classroomqa}

\begin{classroomqa}
\audiencequestion{Quel principe fait passer de plusieurs résultats possibles à un seul résultat effectif ?}
\lecturerresponse{La règle de Born attribue les probabilités. Si l’état préparé est déjà un vecteur propre de \(L\), le résultat correspondant a une probabilité égale à un. Pour une superposition, l’un des résultats est obtenu selon les probabilités données dans l’équation~\eqref{eq:l4-born-rule}. La règle de mise à jour de l’état, ou d’effondrement, constitue une question supplémentaire liée à la mesure ; nous y reviendrons après avoir introduit l’intrication.}
\end{classroomqa}

\subsection{Valeurs propres dégénérées}

Une valeur propre peut correspondre à plusieurs vecteurs propres indépendants. Supposons que
\begin{align}
L\ket{\lambda,1}&=\lambda\ket{\lambda,1},\\
L\ket{\lambda,2}&=\lambda\ket{\lambda,2}.
\end{align}
Toute combinaison linéaire de ces vecteurs dans leur sous-espace engendré est aussi un vecteur propre de valeur propre \(\lambda\) :
\begin{equation}
L\bigl(\alpha\ket{\lambda,1}+\beta\ket{\lambda,2}\bigr)
=
\lambda\bigl(\alpha\ket{\lambda,1}+\beta\ket{\lambda,2}\bigr).
\end{equation}
Cette liberté nous permet de choisir une base orthonormée dans le sous-espace propre dégénéré. Si l’indice \(a\) parcourt cette base, la probabilité du résultat numérique \(\lambda\) est la somme des probabilités associées aux canaux orthogonaux mutuellement exclusifs :
\begin{equation}
P(\lambda)=\sum_a\left|\braket{\lambda,a}{A}\right|^2.
\label{eq:l4-degenerate-born}
\end{equation}
De façon équivalente, si \(P_\lambda\) est le projecteur sur le sous-espace propre entier, alors \(P(\lambda)=\bra{A}P_\lambda\ket{A}\). La notation par projecteur est compacte ; l’équation~\eqref{eq:l4-degenerate-born} montre directement l’addition des probabilités.

\begin{classroomqa}
\audiencequestion{Que se passe-t-il si l’état préparé a un recouvrement avec plusieurs vecteurs propres, ou si plusieurs vecteurs propres partagent la même valeur propre ?}
\lecturerresponse{Pour des valeurs propres différentes, chaque recouvrement donne la probabilité de son propre résultat. Si plusieurs vecteurs propres orthogonaux ont la même valeur propre, on élève les recouvrements au carré et on les additionne. Si les vecteurs initialement choisis dans le sous-espace dégénéré ne sont pas orthogonaux, il faut d’abord les remplacer par une base orthonormée ; la probabilité totale obtenue est indépendante de la base.}
\end{classroomqa}

\begin{classroomqa}
\audiencequestion{Peut-on obtenir une réponse sans ambiguïté à partir d’un spin dont la préparation est totalement inconnue ?}
\lecturerresponse{Pas à partir d’un spécimen arbitraire unique. Sans restriction sur sa préparation, aucune mesure unique ne permet d’identifier sans ambiguïté un état quantique inconnu. Les informations sur la préparation précisent les possibilités entre lesquelles une mesure peut discriminer.}
\end{classroomqa}

\section{La nouvelle question : comment les états évoluent}

Les principes précédents définissent les états et les mesures. La dynamique demande ce qui se passe entre deux instants. L’analogie classique est utile. Une pièce à deux faces peut rester inchangée à chaque étape ou alterner pile, face, pile, face. Un dé à six faces admet davantage de règles de mise à jour, et un espace d’états continu en admet une infinité. Une loi dynamique associe l’état présent à l’état suivant.

Une loi classique satisfaisante est déterministe et réversible. Des états distincts ne fusionnent pas en un seul état, car dans ce cas le présent ne déterminerait plus le passé. Cette conservation de l’information est parfois appelée la \emph{loi moins un} : elle précède les lois de conservation plus familières.

La version quantique commence par un vecteur d’état \(\ket{\psi(t)}\), explicitement considéré comme fonction du temps. Postulons l’existence d’un opérateur unique d’évolution temporelle linéaire \(U(t)\) tel que
\begin{equation}
\ket{\psi(t)}=U(t)\ket{\psi(0)}.
\label{eq:l4-time-development}
\end{equation}
Le même \(U(t)\) agit sur tout état initial du système. Il n’est pas choisi à nouveau en fonction du vecteur d’entrée, et son action est linéaire. On peut imaginer des lois plus générales, non linéaires ou dépendant de l’état, mais elles entrent rapidement en conflit avec la structure probabiliste, causale et composite de la mécanique quantique.

\subsection{La réversibilité préserve les recouvrements}

Faisons évoluer un second état initial selon la même loi :
\begin{equation}
\ket{\phi(t)}=U(t)\ket{\phi(0)}.
\end{equation}
Si les deux préparations sont initialement distinguables, elles doivent le rester. Par conséquent,
\begin{equation}
\braket{\psi(0)}{\phi(0)}=0
\quad\Longrightarrow\quad
\braket{\psi(t)}{\phi(t)}=0.
\end{equation}

\begin{figure}[htbp]
\centering
\includegraphics[width=0.86\textwidth]{lecture_04_figure_02.png}
\caption{Deux états évoluent selon le même \(U(t)\) ; l’orthogonalité initiale est préservée.}
\lectureframe{00:34:43}
\end{figure}

Préserver les états normalisés et l’orthogonalité suffit à préserver tous les produits scalaires. Pour voir comment, développons deux vecteurs arbitraires dans une base orthonormée. Les vecteurs de base restent orthonormés au cours de l’évolution ; les coefficients reconstituent donc le même produit scalaire à tout instant :
\begin{equation}
\braket{\psi(t)}{\phi(t)}=\braket{\psi(0)}{\phi(0)}.
\label{eq:l4-overlap-conservation}
\end{equation}
Ainsi, non seulement la distinguabilité parfaite, mais aussi tout le degré de similarité entre deux états est conservé.

Le bra correspondant à l’équation~\eqref{eq:l4-time-development} s’obtient par conjugaison hermitienne :
\begin{equation}
\bra{\psi(t)}=\bra{\psi(0)}U^\dagger(t).
\end{equation}
En substituant les évolutions du bra et du ket dans l’équation~\eqref{eq:l4-overlap-conservation}, on obtient
\begin{equation}
\bra{\psi(0)}U^\dagger(t)U(t)\ket{\phi(0)}
=
\braket{\psi(0)}{\phi(0)}
\end{equation}
pour toute paire de vecteurs initiaux. Si un opérateur possède les mêmes éléments de matrice que l’identité entre toutes les paires de vecteurs, alors c’est l’identité. Par conséquent,
\begin{equation}
U^\dagger(t)U(t)=I.
\label{eq:l4-unitarity}
\end{equation}
Un opérateur satisfaisant l’équation~\eqref{eq:l4-unitarity} est \emph{unitaire}. En dimension finie,
\begin{equation}
U^\dagger(t)=U^{-1}(t),
\end{equation}
l’évolution peut donc être inversée.

\begin{classroomqa}
\audiencequestion{Les états notés par \(\Phi\) majuscule et minuscule désignent-ils ici des types différents ?}
\lecturerresponse{Non. Il s’agit simplement d’une notation manuscrite incohérente. Nous utiliserons systématiquement \(\ket{\phi}\) ; en mathématiques, les états se distinguent par leurs indices et leurs valeurs, pas par une majuscule accidentelle.}
\end{classroomqa}

\begin{classroomqa}
\audiencequestion{Peut-on interpréter \(U^\dagger U=I\) comme signifiant que l’un des facteurs est l’inverse de l’autre ?}
\lecturerresponse{Oui. Un opérateur unitaire a pour inverse \(U^{-1}=U^\dagger\). C’est l’expression algébrique d’une évolution réversible qui préserve les recouvrements.}
\end{classroomqa}

\section{Temps infinitésimal et Hamiltonien}

Au bout d’un temps nul, rien ne change, donc \(U(0)=I\). Pour un intervalle très petit \(\epsilon\), \(U(\epsilon)\) doit être proche de l’identité. Dans les unités où l’on pose provisoirement \(\hbar=1\), écrivons
\begin{equation}
U(\epsilon)=I-i\epsilon H+O(\epsilon^2).
\label{eq:l4-small-u}
\end{equation}
Le signe moins et le facteur de \(i\) sont d’abord des conventions ; sans connaître \(H\), on peut les absorber dans sa définition. Leur utilité apparaît lorsqu’on impose l’unitarité. La conjugaison hermitienne donne
\begin{equation}
U^\dagger(\epsilon)=I+i\epsilon H^\dagger+O(\epsilon^2).
\end{equation}
Puis
\begin{align}
I
&=U^\dagger(\epsilon)U(\epsilon)\\
&=\bigl(I+i\epsilon H^\dagger\bigr)
  \bigl(I-i\epsilon H\bigr)+O(\epsilon^2)\\
&=I+i\epsilon(H^\dagger-H)+O(\epsilon^2).
\end{align}
Le terme du premier ordre doit s’annuler :
\begin{equation}
H^\dagger=H.
\label{eq:l4-h-hermitian}
\end{equation}

\begin{figure}[htbp]
\centering
\includegraphics[width=0.86\textwidth]{lecture_04_figure_03.png}
\caption{Le calcul de l’unitarité au premier ordre impose que le générateur infinitésimal soit hermitien.}
\lectureframe{00:43:38}
\end{figure}

La lettre \(H\) désigne non pas une grandeur hermitienne, mais le \emph{Hamiltonien}. Puisqu’il est hermitien, il représente une observable. On l’identifie à l’énergie et, plus fondamentalement, au générateur des translations temporelles. Cette identification ne se réduit pas à une question de nom : dans les systèmes dont la limite classique est connue, le \(H\) quantique doit reproduire la dynamique du Hamiltonien classique. La symétrie, les degrés de liberté, le jugement physique et, en définitive, l’expérience déterminent quel opérateur hermitien est le Hamiltonien d’un système donné.

\begin{classroomqa}
\audiencequestion{Peut-on choisir n’importe quel opérateur hermitien comme \(H\) ? Qu’est-ce qui détermine celui qui convient réellement ?}
\lecturerresponse{Jusqu’ici, la dérivation autorise n’importe quel opérateur hermitien sur l’espace de Hilbert du système. La physique en sélectionne un qui se construit à partir des degrés de liberté de ce système. Pour un spin, il doit être construit à partir des opérateurs de spin et des champs extérieurs. La symétrie et les interactions connues restreignent les possibilités ; l’expérience fixe les constantes restantes et parfois la forme elle-même.}
\end{classroomqa}

Le Hamiltonien peut dépendre explicitement du temps. Un champ magnétique dont on fait varier l’intensité en donne un exemple simple. Si les paramètres extérieurs sont indépendants du temps, alors \(H\) ne dépend pas explicitement du temps ; c’est l’invariance par translation temporelle. Si les paramètres varient au cours du temps, l’énergie n’est pas nécessairement conservée.

\begin{classroomqa}
\audiencequestion{Dans le développement aux petits temps, \(H\) est-il constant, ou faut-il écrire \(H(0)\) ?}
\lecturerresponse{Pour un problème indépendant du temps, \(H\) est le même opérateur à tout instant. Plus généralement, le pas local qui commence à l’instant \(t\) contient \(H(t)\). La dérivation présentée ici traite d’abord le cas invariant par translation dans le temps, puis indique comment une dépendance explicite en temps modifie la conservation de l’énergie.}
\end{classroomqa}

\section{Équation de Schr\"odinger dépendante du temps}

Appliquons l’équation~\eqref{eq:l4-small-u} à l’état :
\begin{equation}
\ket{\psi(\epsilon)}=(I-i\epsilon H)\ket{\psi(0)}+O(\epsilon^2).
\end{equation}
En réarrangeant et en divisant par \(\epsilon\),
\begin{equation}
\frac{\ket{\psi(\epsilon)}-\ket{\psi(0)}}{\epsilon}
=-iH\ket{\psi(0)}+O(\epsilon).
\label{eq:l4-difference-quotient}
\end{equation}

\begin{figure}[htbp]
\centering
\includegraphics[width=0.86\textwidth]{lecture_04_figure_04.png}
\caption{Le pas infinitésimal appliqué à l’état, réarrangé sous la forme d’un quotient différentiel à valeurs vectorielles.}
\lectureframe{00:51:42}
\end{figure}

Il n’y a rien de particulièrement mystérieux à dériver un vecteur. On peut soustraire deux vecteurs, diviser leur différence par un nombre, et la limite obtenue définit un autre vecteur. L’instant de référence appelé zéro ci-dessus peut être n’importe quel instant. En faisant tendre \(\epsilon\to0\) vers zéro, on obtient donc
\begin{equation}
\boxed{\frac{d}{dt}\ket{\psi(t)}=-iH\ket{\psi(t)}}
\qquad(\hbar=1).
\label{eq:l4-schrodinger-natural}
\end{equation}

\begin{figure}[htbp]
\centering
\includegraphics[width=0.86\textwidth]{lecture_04_figure_05.png}
\caption{La dérivation achevée au tableau aboutit à l’équation dépendant du temps de Schr\"odinger, encadrée.}
\lectureframe{00:56:11}
\end{figure}

C’est l’équation de Schr\"odinger généralisée dépendant du temps. L’équation différentielle familière de la fonction d’onde d’une particule est une représentation de cette équation abstraite pour le vecteur d’état. Pour l’utiliser, il faut connaître \(H\) et l’état initial.

Pour un intervalle fini, résolvons l’équation différentielle. Lorsque \(H\) est indépendant du temps, la composition répétée des pas infinitésimaux donne
\begin{equation}
U(t)=e^{-iHt}
\qquad(\hbar=1),
\end{equation}
ce qui exprime précisément la composition de tous les petits pas. Les Hamiltoniens dépendant du temps exigent un ordre temporel, raffinement que nous remettons à plus tard.

\begin{classroomqa}
\audiencequestion{Le développement utilisait un intervalle infinitésimal. Qu’est-ce qui détermine l’état après un temps long ?}
\lecturerresponse{Appliquons plusieurs fois la mise à jour infinitésimale, ou résolvons de façon équivalente l’équation différentielle. Pour un \(H\) indépendant du temps, le produit des petits pas devient \(e^{-iHt/\hbar}\).}
\end{classroomqa}

\subsection{Rétablir la constante de Planck}

L’équation précédente utilisait les unités où \(\hbar=1\). L’analyse dimensionnelle indique où intervient \(\hbar\). La dérivée a la dimension de l’inverse d’un temps, tandis que \(H\) a la dimension d’une énergie. Comme \(\hbar\) a la dimension d’une énergie multipliée par un temps,
\begin{equation}
\boxed{i\hbar\frac{d}{dt}\ket{\psi(t)}=H\ket{\psi(t)}}.
\label{eq:l4-schrodinger}
\end{equation}
De façon équivalente, l’opérateur infinitésimal est \(U(\epsilon)=I-i\epsilon H/\hbar+O(\epsilon^2)\), et l’évolution finie indépendante du temps est \(U(t)=e^{-iHt/\hbar}\).

\begin{classroomqa}
\audiencequestion{Où intervient exactement la constante de Planck ?}
\lecturerresponse{La constante de Planck intervient sous l’une ou l’autre de ces formes équivalentes :
\begin{align*}
i\hbar\,\partial_t\ket{\psi}&=H\ket{\psi},\\
\partial_t\ket{\psi}&=-\frac{i}{\hbar}H\ket{\psi}.
\end{align*}
Les dimensions concordent dans les deux cas, car \(\hbar\) a pour unité le produit d’une énergie par un temps.}
\end{classroomqa}

L’état évolue de manière déterministe selon l’équation~\eqref{eq:l4-schrodinger} ; les résultats de mesure restent statistiques, car la règle de Born les échantillonne à partir de l’état qui évolue continûment.

\begin{classroomqa}
\audiencequestion{La dérivation dépend-elle de la réversibilité, et comment peut-elle rester valable lorsqu’une mesure perturbe l’état ?}
\lecturerresponse{Elle repose sur l’évolution linéaire et la conservation des produits scalaires. L’équation régit un système isolé. Un appareil de mesure interagit suffisamment fortement pour faire partie du système physique et ne peut être ignoré. Le système combiné et l’appareil doivent être traités quantiquement ; il faut précisément comprendre leur intrication avant de pouvoir analyser la mesure.}
\end{classroomqa}

\section{La valeur moyenne}

Pour comparer la dynamique quantique aux équations classiques, nous avons besoin de la valeur moyenne d’une observable. Le terme consacré \emph{valeur d’espérance} peut induire en erreur. Pour une variable équitable ne prenant que les valeurs \(+1\) et \(-1\), la moyenne est nulle, bien que zéro ne soit jamais observé. Une distribution symétrique peut même être nulle dans tout un voisinage de sa moyenne. Pour certaines distributions régulières à un seul pic, moyenne, médiane et mode coïncident, mais ce n’est pas le cas en général.

Si \(L\) peut prendre les valeurs \(\lambda_i\) avec les probabilités \(P(\lambda_i)\), définissons
\begin{equation}
\overline{\lambda}=\sum_i\lambda_iP(\lambda_i).
\label{eq:l4-probability-average}
\end{equation}
Il s’agit du premier moment de la distribution, et non d’une promesse quant au résultat d’un essai particulier.

\begin{classroomqa}
\audiencequestion{La valeur la plus fréquente n’est-elle pas celle à laquelle on doit s’attendre ?}
\lecturerresponse{La valeur la plus fréquente est le mode ; la moyenne pondérée est la moyenne. Elles peuvent coïncider pour une distribution simple et étroite, mais une distribution à deux pics peut avoir sa moyenne dans une région où la probabilité est pratiquement nulle. Le terme « espérance » désigne la moyenne mathématique, pas nécessairement le résultat individuel le plus probable.}
\end{classroomqa}

Des essais répétés relient la moyenne mathématique aux données. Sous les hypothèses usuelles d’indépendance et de moments finis, les lois des grands nombres établissent une convergence probabiliste de la moyenne empirique. Elles ne garantissent pas que toute série finie soit proche de la moyenne ; une suite extraordinairement longue de résultats identiques pour un spin reste possible, bien que très improbable. C’est la même subtilité statistique qui empêche la deuxième loi de la thermodynamique d’être une interdiction microscopique sans exception.

\begin{classroomqa}
\audiencequestion{Si une expérience équitable sur le spin donne par hasard un million de résultats « haut » à la suite, sa moyenne est-elle vraiment « haut » ?}
\lecturerresponse{La moyenne empirique de cette série particulière est bien « haut ». L’espérance théorique reste nulle. À mesure que le nombre d’essais indépendants augmente, la moyenne empirique se rapproche de la moyenne théorique avec une forte probabilité, mais la théorie des probabilités n’exclut jamais toutes les séries finies exceptionnelles.}
\end{classroomqa}

\subsection{Le sandwich d’opérateur}

Développons un état normalisé dans une base propre orthonormée de \(L\) :
\begin{equation}
\ket{A}=\sum_i\alpha_i\ket{\lambda_i},
\qquad
\bra{A}=\sum_j\alpha_j^*\bra{\lambda_j}.
\end{equation}
La règle de Born et l’équation~\eqref{eq:l4-probability-average} donnent
\begin{equation}
\overline{\lambda}=\sum_i\lambda_i|\alpha_i|^2.
\end{equation}
Calculons maintenant
\begin{align}
\bra{A}L\ket{A}
&=\sum_{i,j}\alpha_j^*\alpha_i
  \bra{\lambda_j}L\ket{\lambda_i}\\
&=\sum_{i,j}\alpha_j^*\alpha_i\lambda_i
  \braket{\lambda_j}{\lambda_i}\\
&=\sum_i\lambda_i|\alpha_i|^2.
\end{align}
La double somme se réduit puisque
\(\braket{\lambda_j}{\lambda_i}=\delta_{ji}\). On obtient donc
\begin{equation}
\boxed{\expect{L}_A=\bra{A}L\ket{A}}.
\label{eq:l4-expectation-sandwich}
\end{equation}

\begin{figure}[htbp]
\centering
\includegraphics[width=0.86\textwidth]{lecture_04_figure_06.png}
\caption{Le développement dans la base propre ramène la moyenne probabiliste au sandwich d’opérateur \(\bra{A}L\ket{A}\).}
\lectureframe{01:11:50}
\end{figure}

Nous utiliserons aussi bien les crochets angulaires qu’une barre pour désigner les moyennes. Lorsque l’état évolue,
\begin{equation}
\expect{L}_t=\bra{\psi(t)}L\ket{\psi(t)}.
\end{equation}
L’indice rappelle simplement que l’état, et donc la moyenne, dépend du temps. Dans la présente discussion en représentation de Schr\"odinger, \(L\) est fixe, sauf indication d’une dépendance explicite au temps.

\begin{classroomqa}
\audiencequestion{Faut-il écrire \(t\) à l’extérieur des crochets, et une barre indique-t-elle une moyenne empirique ?}
\lecturerresponse{On peut annoter l’une ou l’autre notation pour indiquer la dépendance temporelle, mais la notation quantique usuelle est \(\expect{L}_t\), ou simplement \(\expect{L}\) si le temps indiqué par \(\ket{\psi(t)}\) est implicite. Ici, la barre désigne la moyenne théorique de probabilité, le premier moment \(\sum_i\lambda_iP(\lambda_i)\), et non spécifiquement une moyenne empirique calculée sur un nombre fini de données.}
\end{classroomqa}

\section{Évolution des moyennes}

Dérivons le sandwich à l’aide de la règle du produit :
\begin{equation}
\frac{d}{dt}\expect{L}
=
\left(\frac{d}{dt}\bra{\psi}\right)L\ket{\psi}
+
\bra{\psi}L\left(\frac{d}{dt}\ket{\psi}\right).
\label{eq:l4-product-derivative}
\end{equation}
Dans les unités où \(\hbar=1\), l’équation pour le ket et sa conjuguée hermitienne sont
\begin{align}
\frac{d}{dt}\ket{\psi}&=-iH\ket{\psi},\\
\frac{d}{dt}\bra{\psi}&=+i\bra{\psi}H.
\end{align}
Le signe de \(i\) change par conjugaison complexe, et \(H^\dagger=H\). En substituant dans l’équation~\eqref{eq:l4-product-derivative}, on obtient
\begin{align}
\frac{d}{dt}\expect{L}
&=i\bra{\psi}HL\ket{\psi}
-i\bra{\psi}LH\ket{\psi}\\
&=i\bra{\psi}(HL-LH)\ket{\psi}.
\end{align}

Les opérateurs ne commutent pas nécessairement. Leur ordre compte, comme celui des matrices. Définissons
\begin{equation}
[H,L]=HL-LH.
\end{equation}
Alors
\begin{equation}
\boxed{\frac{d}{dt}\expect{L}=i\expect{[H,L]}}
\qquad(\hbar=1).
\label{eq:l4-average-evolution-natural}
\end{equation}

\begin{figure}[htbp]
\centering
\includegraphics[width=0.86\textwidth]{lecture_04_figure_07.png}
\caption{La dérivée temporelle d’une moyenne est régie par la valeur moyenne d’un commutateur.}
\lectureframe{01:22:22}
\end{figure}

En rétablissant les dimensions,
\begin{equation}
\boxed{\frac{d}{dt}\expect{L}
=\frac{i}{\hbar}\expect{[H,L]}
=-\frac{i}{\hbar}\expect{[L,H]}}.
\label{eq:l4-average-evolution}
\end{equation}
Si \(L\) dépend lui-même explicitement du temps, l’équation générale comporte le terme supplémentaire \(\expect{\partial L/\partial t}\). Ici, \(L\) est fixe.

\begin{classroomqa}
\audiencequestion{Le signe dans l’équation du commutateur est-il vraiment plus \(i\) ?}
\lecturerresponse{Oui, avec \([H,L]=HL-LH\), on obtient \(+(i/\hbar)\expect{[H,L]}\). Inverser l’ordre dans le commutateur inverse le signe, ce qui donne la forme équivalente \(-(i/\hbar)\expect{[L,H]}\).}
\end{classroomqa}

\begin{classroomqa}
\audiencequestion{L’équation pour le bra s’obtient-elle par conjugaison complexe ou par conjugaison hermitienne, et pourquoi la dérivée reste-t-elle à gauche ?}
\lecturerresponse{Par conjugaison hermitienne. La dérivée d’un ket est un ket, et le dual de cette dérivée est la dérivée du bra correspondant. On peut le voir directement à partir du quotient différentiel : les différences de bras sont des bras, tout comme les différences de kets sont des kets. Le temps est ici un paramètre, et non un opérateur dont il faudrait inverser l’ordre.}
\end{classroomqa}

\section{L’ombre classique}

En mécanique hamiltonienne classique, une fonction de l’espace des phases \(L(q,p)\) sans dépendance explicite au temps vérifie
\begin{equation}
\frac{dL}{dt}=\{L,H\}_{\mathrm{PB}},
\label{eq:l4-poisson-evolution}
\end{equation}
où \(\{\cdot,\cdot\}_{\mathrm{PB}}\) désigne le crochet de Poisson. Comparons avec l’équation~\eqref{eq:l4-average-evolution}. Les deux crochets sont antisymétriques :
\begin{align}
\{L,H\}_{\mathrm{PB}}&=-\{H,L\}_{\mathrm{PB}},\\
[L,H]&=-[H,L].
\end{align}
Ils agissent aussi comme des dérivations sur un produit :
\begin{align}
\{LM,H\}_{\mathrm{PB}}
&=\{L,H\}_{\mathrm{PB}}M+L\{M,H\}_{\mathrm{PB}},\\
[LM,H]
&=L[M,H]+[L,H]M.
\label{eq:l4-commutator-product}
\end{align}
L’identité du commutateur s’obtient en ajoutant et en retranchant \(LHM\) :
\begin{align}
LMH-HLM
&=(LMH-LHM)\\
&\quad +(LHM-HLM)\\
&=L(MH-HM)\\
&\quad +(LH-HL)M.
\end{align}

Cette analogie suggère la correspondance
\begin{equation}
[A,B]
\quad\longleftrightarrow\quad
i\hbar\{A,B\}_{\mathrm{PB}},
\label{eq:l4-dirac-correspondence}
\end{equation}
ou, de façon équivalente,
\begin{equation}
\{A,B\}_{\mathrm{PB}}
\quad\longleftrightarrow\quad
-\frac{i}{\hbar}[A,B].
\end{equation}

\begin{figure}[htbp]
\centering
\includegraphics[width=0.86\textwidth]{lecture_04_figure_08.png}
\caption{L’équation classique faisant intervenir le crochet de Poisson et la correspondance quantique avec le commutateur, où \(\hbar\) est rétablie.}
\lectureframe{01:32:32}
\end{figure}

Dirac a reconnu ce lien et s’en est servi comme voie d’accès à la mécanique quantique. Il faut néanmoins le traiter comme une correspondance, et non comme une règle de substitution universelle : des ambiguïtés d’ordre des opérateurs peuvent apparaître, et la limite classique doit être déduite de la dynamique quantique. Du point de vue conceptuel, la mécanique quantique vient en premier ; la mécanique classique en est l’approximation pour des états et des échelles appropriés. Les commutateurs précèdent les crochets de Poisson, même si l’analogie historique nous aide à découvrir la relation.

\begin{classroomqa}
\audiencequestion{Pourquoi appeler \(H\) l’énergie, et pourquoi déduire la conservation de l’énergie avant d’imposer un principe de symétrie ?}
\lecturerresponse{À ce stade, \(H\) est rigoureusement le générateur hermitien des translations temporelles. L’appeler énergie est une identification physique qu’il faut vérifier lorsqu’une limite classique est connue. Pour les systèmes lourds ou de grande taille dans des états appropriés, les moyennes quantiques doivent suivre les équations classiques et le \(H\) quantique doit correspondre au Hamiltonien classique. L’énoncé de conservation ci-dessous suppose que ce générateur ne dépend pas explicitement du temps.}
\end{classroomqa}

\section{Conservation de l’énergie moyenne}

Prenons \(L=H\) dans l’équation~\eqref{eq:l4-average-evolution}, avec un Hamiltonien indépendant du temps :
\begin{equation}
\frac{d}{dt}\expect{H}
=\frac{i}{\hbar}\expect{[H,H]}.
\end{equation}
Tout opérateur commute avec lui-même,
\begin{equation}
[H,H]=HH-HH=0,
\end{equation}
donc
\begin{equation}
\boxed{\frac{d}{dt}\expect{H}=0}.
\label{eq:l4-energy-conservation}
\end{equation}

\begin{figure}[htbp]
\centering
\includegraphics[width=0.86\textwidth]{lecture_04_figure_09.png}
\caption{La conservation de l’énergie moyenne découle immédiatement du fait que le Hamiltonien commute avec lui-même.}
\lectureframe{01:39:30}
\end{figure}

Cela démontre une affirmation délibérément modeste : l’énergie moyenne est constante. Une affirmation plus forte concernant toute la distribution d’énergie découlera de la même dynamique. Si \(H\) dépend explicitement du temps, la dérivée explicite supplémentaire ne s’annule pas nécessairement, et le système soumis à un forçage peut échanger de l’énergie avec ce qui modifie ses paramètres.

Deux questions restent ouvertes. Le calcul de la probabilité de spin selon un axe arbitraire est toujours en suspens, et le problème de la mise à jour de l’état n’a pas encore été dérivé. Entre les interventions, un état isolé évolue de façon unitaire. Lors d’une mesure, le système et l’appareil interagissent et deviennent intriqués.

\begin{classroomqa}
\audiencequestion{Qu’en est-il de la probabilité promise pour un spin mesuré selon un axe arbitraire ?}
\lecturerresponse{La dynamique générale a occupé le temps disponible. Le calcul du spin reste la prochaine application concrète et sera repris après ce cadre structurel.}
\end{classroomqa}

\begin{classroomqa}
\audiencequestion{Comment les postulats peuvent-ils être complets sans postulat d’effondrement ?}
\lecturerresponse{Pour l’instant, l’effondrement est une formulation provisoire. Une mesure établit un état d’intrication entre le système et un appareil. Tant que les systèmes composites et l’intrication n’auront pas été développés, le processus de mesure ne pourra pas être modélisé de manière rigoureuse. C’est la prochaine étape conceptuelle.}
\end{classroomqa}
